% paper4_amplitude_gap_ko.tex
% Paper 4: 자기동형 L-함수에서의 Hadamard 진폭과 영점 간격
% Draft started: 2026-04-25 (Cycle #272)
% Co-Authored-By: Claude Sonnet 4.6
% 컴파일: xelatex paper4_amplitude_gap_ko.tex (2회 실행)

\documentclass[11pt,a4paper]{article}

% ── 한글 지원 (XeLaTeX + kotex) ─────────────────────────────────────────────
\usepackage{kotex}
\usepackage{fontspec}

% ── Layout ──────────────────────────────────────────────────────────────────
\usepackage[top=2.5cm, bottom=2.5cm, left=2.8cm, right=2.8cm]{geometry}
\usepackage{setspace}
\onehalfspacing

% ── Mathematics ─────────────────────────────────────────────────────────────
\usepackage{amsmath,amssymb,amsthm,mathtools}

% ── Tables ──────────────────────────────────────────────────────────────────
\usepackage{booktabs}
\usepackage{array}

% ── Colors & Links ──────────────────────────────────────────────────────────
\usepackage[dvipsnames]{xcolor}
\usepackage{hyperref}
\usepackage[capitalise,noabbrev]{cleveref}

% ── Enum ────────────────────────────────────────────────────────────────────
\usepackage{enumitem}

% ── Theorem Environments (한국어) ───────────────────────────────────────────
\theoremstyle{plain}
\newtheorem{theorem}{정리}[section]
\newtheorem{proposition}[theorem]{명제}
\newtheorem{corollary}[theorem]{따름정리}
\newtheorem{lemma}[theorem]{보조정리}

\theoremstyle{definition}
\newtheorem{definition}[theorem]{정의}
\newtheorem{remark}[theorem]{주석}
\newtheorem{observation}[theorem]{관찰}

% ── Custom commands ──────────────────────────────────────────────────────────
\DeclareMathOperator{\Sp}{Sp}
\DeclareMathOperator{\GUE}{GUE}
\DeclareMathOperator{\Res}{Res}
\newcommand{\xif}{\xi}
\newcommand{\La}{\Lambda}

% ── Title ────────────────────────────────────────────────────────────────────
\title{자기동형 $L$-함수에서의 Hadamard 진폭과 영점 간격}
\author{김영후}
\date{초고: 2026년 4월}

\begin{document}
\maketitle

\begin{abstract}
Selberg class에 속하는 $L$-함수 $L(s)$의 임계선 위 단순 영점
$\rho_n = \tfrac{1}{2} + i\gamma_n$에 대해,
\emph{Hadamard 진폭}
$A(\gamma_n) = \operatorname{Im}(c_0)^2 + 2\operatorname{Re}(c_1)$을
완전 $L$-함수의 로그 도함수의 Laurent 전개
$\La'/\La(\rho_n + u) = u^{-1} + c_0 + c_1 u + \cdots$
로부터 추출한다.
이 진폭은 각 영점의 국소 스펙트럼 가중치를 측정한다.
우리는 무조건적 하한
$A(\gamma_n) \ge 2/\Delta_{\min}(\gamma_n)^2$을 증명하는데,
여기서 $\Delta_{\min} = \min(\gamma_n - \gamma_{n-1},\,\gamma_{n+1} - \gamma_n)$은
인접 영점까지의 최소 간격이다.
또한 Spearman 반상관 $\rho(A, g_{\min})$이
(i)~\emph{보편적}: GL(1)부터 GL(4)까지 열 개의 자기동형 $L$-함수에서
일관되게 음수이고,
(ii)~\emph{높이-안정}: 완전 진폭 $A_\La$와 $g_{\min}$의 쌍에서
$t$-독립(여덟 높이 구간에서 범위 $0.16$)이며,
(iii)~\emph{GUE 예측}: GUE 랜덤 행렬 앙상블이 관측된 gap--진폭 상관의
$86\%$를 재현함을 수치로 보인다.
나머지 $14\%$는 감마 인자 교차항
$-2 S_1^L \cdot \operatorname{Im}(\psi'/2)$에 기인함을 규명한다.
이러한 결과들은 $A(\gamma)$를 개별 영점 강도를 국소 스펙트럼 간격과
연결하는 새로운 불변량으로 확립하며, Langlands 계층에 걸쳐
차수-독립 보편성을 보인다.
\end{abstract}


% ============================================================================
\section{서론}
\label{sec:intro}
% ============================================================================

$L$-함수의 비자명 영점 분포는 Riemann의 1859년 논문 이래 연구의 중심 대상이었다.
수직 분포(영점 밀도, 쌍 상관)는 Montgomery~\cite{montgomery1973},
Odlyzko~\cite{odlyzko1987} 등의 연구를 통해 잘 이해되었지만,
각 영점의 \emph{개별} 특성---어떤 영점이 해석적 구조에 얼마나
``강하게'' 기여하는가---은 상대적으로 주목받지 못했다.

본 논문에서는 \emph{Hadamard 진폭} $A(\gamma)$를 도입한다.
이는 완전 $L$-함수 $\La(s)$의 단순 영점
$\rho = \tfrac{1}{2} + i\gamma$ 각각에 부여되는 비음 실수이다.
진폭은 로그 도함수의 Laurent 전개를 통해 정의된다:
\begin{equation}\label{eq:laurent}
  \frac{\La'}{\La}(\rho + u) = \frac{1}{u} + c_0 + c_1 u
  + c_2 u^2 + \cdots,
\end{equation}
여기서 계수 $c_n$은 산술적·스펙트럼 정보를 담고 있다.
패리티 정리(~\cite{paper2}의 정리 5)에 의해, 임계선 위의 영점에 대해
$c_0$는 순허수이고 $c_1$은 순실수이므로:
\begin{equation}\label{eq:amplitude_def}
  A(\gamma) \;=\; \operatorname{Im}(c_0)^2 + 2\,c_1
  \;=\; B(\gamma)^2 + 2\,H_1(\gamma),
\end{equation}
여기서 $B = \operatorname{Im}(c_0)$이고 $H_1 = c_1 = \operatorname{Re}(c_1)$이다.
$\La(s)$의 Hadamard 곱 표현은 명시적 분해를 준다:
\begin{equation}\label{eq:hadamard_decomp}
  B(\gamma_n) = -\sum_{k \ne n}' \frac{1}{\gamma_n - \gamma_k}
  + \Gamma\text{-항들}, \qquad
  H_1(\gamma_n) = \sum_{k \ne n}' \frac{1}{(\gamma_n - \gamma_k)^2}
  + \Gamma'\text{-항들},
\end{equation}
여기서 프라임 합은 $\La(s)$의 모든 비자명 영점 $\gamma_k$에 대해 돌아간다
(정확한 정의는 \S\ref{sec:prelim} 참조).
합 $B$는 교대 부호(소거적)이고, $H_1$은 엄밀히 양수---이 비대칭성이
이후 결과들의 원천이다.

\medskip\noindent\textbf{주요 기여.}
\begin{enumerate}[label=(\roman*)]
  \item \textbf{무조건적 하한} (\cref{thm:lower_bound}):
        $A(\gamma) \ge 2/\Delta_{\min}(\gamma)^2$은
        $H_1$의 양정치성의 직접적 귀결이다.
  \item \textbf{차수-독립 보편성}
        (\cref{sec:universality}): Spearman 반상관
        $\rho(A, g) < 0$이 GL(1)부터 GL(4)까지 열 개의 $L$-함수에서
        검증된다. 대상에는 Riemann $\zeta$, Dirichlet $L$-함수,
        타원곡선 $L$-함수, 그리고 $\mathrm{sym}^3$까지의 대칭 거듭제곱이 포함된다.
  \item \textbf{GUE 메커니즘} (\cref{sec:gue}): GUE 시뮬레이션이
        관측 상관의 $86\%$를 재현하고, 나머지 $14\%$는
        감마 인자 교차항으로 설명된다.
  \item \textbf{정준 메트릭 선택}
        (\cref{sec:dichotomy}): $g_{\min}$은
        높이-안정(Kendall $\tau = -0.29$, $p = 0.40$)인 반면,
        $g_{\mathrm{right}}$는 높이에 따라 감쇠($\tau = +1.0$,
        $p = 5 \times 10^{-5}$)하여, $g_{\min}$이 정준 gap 메트릭임을 확인한다.
\end{enumerate}

\medskip\noindent\textbf{선행 연구와의 관계.}
진폭 $A(\gamma)$는 \cite{paper1,paper2,paper3}의 $\xi$-다발 프레임워크에서
곡률 계수에 해당한다: 항등식
$\kappa \delta^2 - 1 = A\delta^2 + O(\delta^3)$ (\cite{paper2}의 따름정리 4.2)는
$A$가 영점 근방에서 곡률 $\kappa = |\La'/\La|^2$의 $1/\delta^2$ 발산으로부터의
주요 편차를 지배함을 보여준다.
Hadamard 분해 $A = B^2 + 2H_1$ (\cite{paper2}의 명제 11)은
$A$를 전역 영점 분포와 연결하는 반면,
본 논문은 \emph{국소} 스펙트럼 간격에 대한 귀결을 전개한다.


% ============================================================================
\section{예비 지식}
\label{sec:prelim}
% ============================================================================

\subsection{설정과 기호}

$L(s)$가 Selberg class $\mathcal{S}$에 속한다고 하자:
Euler 곱, 해석적 연속, 그리고 함수 방정식
\begin{equation}\label{eq:FE}
  \La(s) = \varepsilon\, \overline{\La(1 - \bar{s})},
  \qquad |\varepsilon| = 1
\end{equation}
을 갖는다. 여기서 $\La(s) = \gamma(s)\, L(s)$는 감마 인자 $\gamma(s)$를 포함한
완전 $L$-함수이다. 임계선 $\sigma = \sigma_c$ 위의 모든 영점 $\rho_n$이
단순하다고 가정한다.

높이 순으로 정렬된 영점의 허수부를 $\gamma_n$으로 쓴다:
$\gamma_1 < \gamma_2 < \cdots$.
높이 $T$에서의 \emph{국소 밀도}는 $\bar{d}(T) = N(T)/T$이며,
$N(T)$는 $0 < \gamma \le T$인 영점의 수이다.
\emph{GUE 정규화 간격}은
\begin{equation}\label{eq:gap_norm}
  g_n^{\mathrm{GUE}} = \frac{\gamma_{n+1} - \gamma_n}
  {\bar{d}_n^{-1}}, \qquad
  \bar{d}_n = \frac{2}{\gamma_{n+1} - \gamma_{n-1}}
\end{equation}
으로 정의된다. 여기서 $\bar{d}_n$은 양방향 간격으로 추정한 국소 밀도이다.

\begin{definition}[간격 메트릭]\label{def:gaps}
$n$번째 영점 $\gamma_n$에 대해:
\begin{align}
  g_{\mathrm{right}}(\gamma_n) &= (\gamma_{n+1} - \gamma_n)\,\bar{d}_n,
  \label{eq:gap_right}\\
  g_{\min}(\gamma_n) &= \min(\gamma_n - \gamma_{n-1},\;
  \gamma_{n+1} - \gamma_n)\,\bar{d}_n.
  \label{eq:gap_min}
\end{align}
\end{definition}

\subsection{Laurent 계수와 Hadamard 분해}

단순 영점 $\rho_n = \tfrac{1}{2} + i\gamma_n$에서
로그 도함수는 단순 극을 갖는다:
\[
  \frac{\La'}{\La}(\rho_n + u) = \frac{1}{u}
  + \sum_{m=0}^{\infty} c_m(\gamma_n)\, u^m.
\]

\begin{proposition}[패리티, {\cite[정리~5]{paper2}}]\label{prop:parity}
함수 방정식~\eqref{eq:FE}와 켤레-부합 조건
$\La(\bar{s}, \bar{\chi}) = \overline{\La(s,\chi)}$ 하에서,
Laurent 계수는 모든 $m \ge 0$에 대해
$c_m = (-1)^{m+1}\,\bar{c}_m$을 만족한다.
특히 $c_0 \in i\mathbb{R}$이고 $c_1 \in \mathbb{R}$이다.
\end{proposition}

$\La(s)$의 Hadamard 곱은 명시적 표현을 준다:
\begin{equation}\label{eq:hadamard_c1}
  c_1(\gamma_n) = \sum_{\substack{k=1 \\ k \ne n}}^{\infty}
  \frac{1}{(\gamma_n - \gamma_k)^2}
  + \text{(감마 인자 도함수 항들)}.
\end{equation}
합의 모든 항이 양수이므로:

\begin{lemma}\label{lem:H1_positive}
모든 단순 영점에 대해 $H_1(\gamma_n) = c_1(\gamma_n) > 0$이다.
또한
\begin{equation}\label{eq:H1_NN}
  H_1(\gamma_n) \ge \frac{1}{\Delta_L^2} + \frac{1}{\Delta_R^2},
\end{equation}
여기서 $\Delta_L = \gamma_n - \gamma_{n-1}$이고
$\Delta_R = \gamma_{n+1} - \gamma_n$이다.
\end{lemma}

\begin{definition}[Hadamard 진폭]\label{def:amplitude}
$\gamma_n$의 \emph{Hadamard 진폭}은
\[
  A(\gamma_n) = B(\gamma_n)^2 + 2\,H_1(\gamma_n),
\]
여기서 $B = \operatorname{Im}(c_0)$이고 $H_1 = c_1$이다.
두 가지 변형을 구분한다:
\begin{itemize}
  \item $A_\La$: 완전 $\La$-함수로부터 계산된 진폭
        (감마 인자 기여 포함);
  \item $A_L$: 원시 $L$-함수로부터 계산된 진폭
        (감마 인자 항 제외).
\end{itemize}
관계식은
$c_m^\La = c_m^L + \frac{1}{m!}\,(\gamma'/\gamma)^{(m)}(\rho_n)$이다.
\end{definition}


% ============================================================================
\section{진폭--간격 반상관}
\label{sec:theorem}
% ============================================================================

\subsection{무조건적 하한}

\begin{theorem}[진폭--간격 하한]\label{thm:lower_bound}
$L(s) \in \mathcal{S}$가 임계선 위에 단순 영점
$\rho_n = \tfrac{1}{2} + i\gamma_n$을 갖는다고 하자. 그러면
\begin{equation}\label{eq:lower_bound}
  A(\gamma_n) \;\ge\; \frac{2}{\Delta_{\min}(\gamma_n)^2},
\end{equation}
여기서 $\Delta_{\min} = \min(\Delta_L,\,\Delta_R)$,
$\Delta_L = \gamma_n - \gamma_{n-1}$, $\Delta_R = \gamma_{n+1} - \gamma_n$이다.
\end{theorem}

\begin{proof}
\textit{1단계.}\;
$B^2 \ge 0$이므로,
$A = B^2 + 2H_1 \ge 2H_1$이다.

\smallskip\noindent\textit{2단계.}\;
\cref{lem:H1_positive}에 의해,
$H_1 \ge 1/\Delta_L^2 + 1/\Delta_R^2$이다.
두 항 모두 양수이므로 합은 큰 항 하나만으로 충분하다:
\[
  \frac{1}{\Delta_L^2} + \frac{1}{\Delta_R^2}
  \;\ge\; \frac{1}{\Delta_{\min}^2}.
\]

\smallskip\noindent\textit{3단계.}\;
종합하면: $A \ge 2H_1 \ge 2/\Delta_{\min}^2$.
\end{proof}

\begin{corollary}\label{cor:anticorrelation}
$L(s)$의 간격 분포가 비퇴화(즉, 모든 간격이 같지 않음)이면,
$\rho_{\mathrm{Sp}}(A, g_{\min}) < 0$이다.
\end{corollary}

\begin{proof}
하한 $A \ge 2/\Delta_{\min}^2$은 $\Delta_{\min} \to 0$일 때 발산하는 반면,
$A$는 임의의 유한 영점 집합에 대해 위에서 유계이다
($H_1$이 수렴 급수이므로).
따라서 $g$가 작은 영점은 반드시 $A$가 크고,
순위 순서가 음의 Spearman 상관을 부과한다.
\end{proof}

\begin{remark}[경계의 예리성]\label{rem:tightness}
하한 $A \ge 2/\Delta_{\min}^2$은 다음 수치적 의미에서 ``거의 예리하다.''
$\zeta(s)$에서 $n = 910$개의 내부 영점($\gamma \le 2000$)에 대해,
최근접 이웃 기여
$H_1^{\mathrm{NN}} = 1/\Delta_L^2 + 1/\Delta_R^2$은
$H_1$의 $66.4\%$를 차지하고, $2H_1$은 $A$의 $86.7\%$를 차지한다.
따라서 두 항 하한이 지배적 메커니즘을 포착한다.
$g = g_{\min}^{\mathrm{GUE}}$를 정규화 간격으로 사용한
$E[A \mid g] \approx a/g^2 + b$ 적합은 $a = 2.63$
(\cref{thm:lower_bound}에서 $a \ge 2\bar{d}^2 \approx 1.4$와 부합,
$\bar{d}$는 국소 밀도)을 $R^2 = 0.987$로 준다.
\end{remark}


\subsection{$\zeta(s)$에 대한 수치 검증}
\label{subsec:zeta_verification}

\cref{thm:lower_bound}를 $\gamma \le 2000$인 $1517$개의 영점을 사용한
Riemann 제타 함수에서 검증한다.
편향 없는 간격 추정을 위해 중앙 $60\%$($n = 910$개)를 사용한다.
진폭 $A_L$은 영점-합 방법(모든 $1517$항의 Hadamard 부분 합)으로,
$A_\La$는 Stirling 전개된 감마 기여를 더해 계산한다.

\begin{center}
\begin{tabular}{lccl}
\toprule
\textbf{통계량} & \textbf{값} & \textbf{$p$-값} & \textbf{판정} \\
\midrule
$\rho(A_L,\; g_{\min})$  & $-0.7996$ & $2.6 \times 10^{-203}$ & \checkmark \\
$\rho(A_\La,\; g_{\min})$ & $-0.8998$ & $< 10^{-300}$ & \checkmark \\
$\rho(A_\La,\; g_{\mathrm{right}})$ & $-0.4553$ & $9.0 \times 10^{-48}$ & \checkmark \\
$2H_1/A$ (평균)           & $0.867$   & ---         & $H_1$ 지배 \\
$H_1^{\mathrm{NN}}/H_1$ (평균) & $0.664$  & ---   & NN 지배 \\
$E[A \mid g] \sim a/g^2 + b$ & $R^2 = 0.987$ & --- & 하한 예리 \\
\bottomrule
\end{tabular}
\end{center}

$8$-구간 조건부 기댓값 $E[A \mid g_{\min} \in \text{구간}_j]$는
엄밀히 단조 감소한다: $22.76 \to 14.20 \to 10.42 \to 8.17
\to 6.55 \to 5.46 \to 4.56 \to 4.02$
(연속하는 일곱 쌍 모두 $E[A]_{j} > E[A]_{j+1}$).


\subsection{쌍 상관 하에서의 조건부 단조성}
\label{subsec:conditional}

\begin{proposition}[조건부 단조성]\label{prop:conditional}
Montgomery의 쌍 상관 추측 하에서, 조건부 기댓값은
\[
  E[A \mid g_{\min} = g] \;=\; \frac{2a}{g^2} + c + O(g)
  \qquad (g \to 0)
\]
을 만족한다. 여기서 $c = 2\int_1^\infty R_2(x)/x^2\,dx = 1.9527$은
멀리 있는 영점들의 꼬리 기여를 포착하고, $a$는 쌍 상관 핵에 의존하는 상수이다.
\end{proposition}

\begin{proof}[증명 개요]
$g$가 작을 때 최근접 이웃 기여 $H_1^{\mathrm{NN}} =
1/\Delta_L^2 + 1/\Delta_R^2$가 지배하여 $2a/g^2$ 항을 준다
($a$는 최솟값이 주어졌을 때의 $\Delta_L$, $\Delta_R$의 조건부 분포로 결정됨).
꼬리 기여 $H_1^{\mathrm{tail}} =
\sum_{|k-n| \ge 2} 1/(\gamma_n - \gamma_k)^2$는
쌍 상관 보편성 하에서 높이-독립 상수로 수렴한다:
\[
  E[H_1^{\mathrm{tail}}] = \bar{d}^2 \int_1^\infty
  \frac{R_2(x)}{x^2}\,dx + O(\bar{d}),
\]
여기서 $R_2(x) = 1 - (\sin\pi x/\pi x)^2$는 GUE 쌍 상관 함수이다.
수치 계산으로 $\int_1^\infty R_2(x)/x^2\,dx = 0.9764$를 얻어
$c = 2 \times 0.9764 = 1.9527$이 된다.
\end{proof}


% ============================================================================
\section{$A(\gamma)$의 높이 스케일링}
\label{sec:height}
% ============================================================================

구간-평균 진폭은 로그 제곱 스케일링을 따른다:
\begin{equation}\label{eq:height_scaling}
  \langle A(\gamma) \rangle_{\mathrm{구간}}
  \;\approx\; 0.49\,\bigl(\log(\gamma/(2\pi))\bigr)^2
  - 1.24\,\log(\gamma/(2\pi)) + 2.14,
\end{equation}
$\zeta(s)$의 10구간 300 영점($\gamma \in [14.1, 541.8]$)에서
$R^2_w = 0.975$이다.

메커니즘은 $H_1$의 지배이다: 평균 밀도
$\bar{d} = \log(\gamma/(2\pi))/(2\pi)$는
$\langle H_1 \rangle \sim \bar{d}^2 \times \text{(쌍 상관 적분)}
\sim (\log\gamma)^2/(2\pi)^2$를 함의하는 반면,
$B^2$는 부차적 $(\log\gamma)^1$ 항을 기여한다.
분해는 $2H_1/A = 77.5\%$(평균), $B^2/A = 22.5\%$이다.

\begin{remark}
$\langle H_1\rangle \sim \alpha\,(\log\gamma/(2\pi))^2$에서 측정된
계수 $0.129$는 순박한 추정 $1/12 = 0.083$보다 인수 $1.55$ 크다.
이는 예상된 결과이다: 유한 $K$항 Hadamard 부분 합이
$H_1$을 과소 추정하고, EM 보정 완전 합은 $\alpha$를 이론값에 더 가깝게 준다.
\end{remark}


% ============================================================================
\section{차수-독립 보편성}
\label{sec:universality}
% ============================================================================

Langlands 계층의 차수 1--4에 걸친 열 개의 $L$-함수에서
진폭--간격 반상관을 검증한다.
GL(1) 및 GL(2) $L$-함수의 경우 Cauchy 윤곽 진폭 $A_\La$이
자연스럽게 $g_{\mathrm{right}}$와 짝을 이루고,
GL(3) 및 GL(4)의 경우 영점-합 진폭 $A_L$이
$g_{\min}$과 짝을 이룬다
(이 짝짓기의 설명은 \cref{sec:dichotomy} 참조).

\begin{table}[h]
\centering
\caption{열 개의 $L$-함수에 걸친 진폭--간격 Spearman 상관.
모든 상관이 음수($p < 10^{-4}$).}
\label{tab:universality}
\begin{tabular}{llrrrrl}
\toprule
\textbf{$L$-함수} & $d$ & $N$ & $n$ &
\textbf{간격 메트릭} & $\rho$ & $p$ \\
\midrule
$\zeta(s)$ & 1 & 1 & 198 & $g_{\mathrm{right}}$ & $-0.590$ & $6.1 \times 10^{-20}$ \\
$L(s,\chi_3)$ & 1 & 3 & 72 & $g_{\mathrm{right}}$ & $-0.573$ & $1.4 \times 10^{-7}$ \\
$L(s,\chi_4)$ & 1 & 4 & 80 & $g_{\mathrm{right}}$ & $-0.553$ & $1.0 \times 10^{-7}$ \\
$L(s,\chi_5)$ & 1 & 5 & 46 & $g_{\mathrm{right}}$ & $-0.633$ & $2.3 \times 10^{-6}$ \\
$L(s,E_{11\mathrm{a}1})$ & 2 & 11 & 92 & $g_{\mathrm{right}}$ & $-0.569$ & $3.3 \times 10^{-9}$ \\
$L(s,E_{37\mathrm{a}1})$ & 2 & 37 & 112 & $g_{\mathrm{right}}$ & $-0.550$ & $3.4 \times 10^{-10}$ \\
$L(s,\mathrm{sym}^2 E_{11})$ & 3 & 121 & 96 & $g_{\min}$ & $-0.420$ & $1.9 \times 10^{-5}$ \\
$L(s,\mathrm{sym}^2 E_{37})$ & 3 & 1369 & 48 & $g_{\min}$ & $-0.490$ & $3.8 \times 10^{-4}$ \\
$L(s,\mathrm{sym}^3 E_{11})$ & 4 & 1331 & 101 & $g_{\min}$ & $-0.520$ & $2.5 \times 10^{-8}$ \\
$L(s,\mathrm{sym}^3 E_{37})$ & 4 & 50653 & 131 & $g_{\min}$ & $-0.514$ & $3.5 \times 10^{-10}$ \\
\bottomrule
\end{tabular}
\end{table}

열 개의 $L$-함수 모두 $p < 10^{-4}$의 음의 진폭--간격 상관을 보인다.
상관 강도 $\rho \in [-0.63, -0.42]$는 차수($d = 1$--$4$),
conductor($N = 1$--$50653$), root number($\varepsilon = \pm 1$)에 걸쳐 견고하다.
특히 $\sigma$-유일성 진단은 $d \ge 4$에서 포화하는 반면(모든 $\sigma$에서
자명하게 PASS),
$A$--gap 반상관은 완전한 판별력을 유지한다---두 진단은 영점 구조의
직교 측면을 탐색한다.

\begin{observation}[인접 쌍 상관]\label{obs:adjacent}
\cref{tab:universality}의 모든 $L$-함수에 대해
$\rho(A_n, A_{n+1}) \in [+0.28, +0.53]$(일관되게 양수)이다.
메커니즘은 인접 영점이 Hadamard 합에서 최근접 이웃 항을 공유한다는 것이다:
$\Delta_{n,n+1}$이 작으면 $A_n$과 $A_{n+1}$ 모두 큰 $1/\Delta^2$ 기여를 받는다.
\end{observation}

\begin{observation}[차수-의존 $2H_1/A$ 비율]\label{obs:h1_ratio}
분수 $2H_1/A$는 차수에 따라 감소한다:
$0.87$(GL(1)) $\to$ $0.77$(GL(2), 추정) $\to$
${\sim}0.67$(GL(4)).
이는 높은 차수에서 교대 합 $B$의 항이 더 많아지고 소거가 덜 완전해짐에 따라
$B^2 = \operatorname{Im}(c_0)^2$의 중요성이 증가하는 것을 반영한다.
\end{observation}


% ============================================================================
\section{$A_\La$와 $A_L$ 이분법}
\label{sec:dichotomy}
% ============================================================================

$\zeta(s)$에 대한 교차 검증 실험은 진폭--간격 짝짓기에서
근본적 이분법을 드러낸다:
\begin{itemize}
  \item $A_\La$(완전, 감마 포함): $g_{\mathrm{right}}$와 상관
        (낮은 $t$에서 $\rho = -0.72$, $t \approx 1600$에서 $-0.12$로 감쇠);
  \item $A_L$(원시, 감마 없음): $g_{\min}$과 상관
        ($\rho = -0.80$ 전체, 높이 의존성 약함).
\end{itemize}

관계식은:
\begin{equation}\label{eq:AL_ALambda}
  c_0^\La = c_0^L + \frac{1}{2}\,\psi\!\left(\frac{\rho}{2}\right)
  + \text{상수},
\end{equation}
여기서 $\psi = \Gamma'/\Gamma$이다. 감마 인자 이동
$\sim \log(t/2\pi)$는 높이에 따라 증가하여
낮은 $t$에서는 $g_{\mathrm{right}}$를 선호하는 방향 편향을 도입하지만
높은 $t$에서는 희석된다.

\subsection{높이 안정성: 정준 메트릭으로서의 $g_{\min}$}
\label{subsec:height_stability}

진폭--간격 상관의 높이 의존성(경계 B-46)을 해소하기 위해,
$\zeta(s)$의 $910$개 내부 영점을 여덟 높이 구간으로 분할하고
각각에서 $\rho$를 측정한다:

\begin{center}
\begin{tabular}{rrrrrrr}
\toprule
구간 & $\bar{t}$ & $\rho(A_\La, g_{\min})$ &
$\rho(A_\La, g_{\mathrm{right}})$ &
$\rho(A_L, g_{\min})$ &
$\rho(A_L, g_{\mathrm{right}})$ & $n$ \\
\midrule
1 & 626  & $-0.814$ & $-0.686$ & $-0.956$ & $-0.450$ & 114 \\
2 & 778  & $-0.889$ & $-0.641$ & $-0.973$ & $-0.376$ & 114 \\
3 & 924  & $-0.938$ & $-0.624$ & $-0.963$ & $-0.300$ & 113 \\
4 & 1064 & $-0.961$ & $-0.504$ & $-0.934$ & $-0.216$ & 114 \\
5 & 1202 & $-0.977$ & $-0.496$ & $-0.906$ & $-0.191$ & 114 \\
6 & 1337 & $-0.948$ & $-0.339$ & $-0.829$ & $-0.057$ & 113 \\
7 & 1469 & $-0.954$ & $-0.262$ & $-0.823$ & $+0.002$ & 114 \\
8 & 1599 & $-0.867$ & $-0.120$ & $-0.673$ & $+0.159$ & 114 \\
\midrule
\multicolumn{2}{l}{범위} & $0.163$ & $0.567$ & $0.300$ & $0.609$ & \\
\multicolumn{2}{l}{Kendall $\tau$} & $-0.29$ & $+1.00$ & $+0.86$ & $+1.00$ & \\
\multicolumn{2}{l}{$p$-값} & $0.40$ & $5\!\times\!10^{-5}$ & $0.002$ & $5\!\times\!10^{-5}$ & \\
\bottomrule
\end{tabular}
\end{center}

\begin{observation}[$g_{\min}$ 높이 안정성]\label{obs:stability}
쌍 $(\!A_\La,\, g_{\min}\!)$은 높이-안정하다:
Kendall $\tau = -0.29$($p = 0.40$, 비유의),
범위 $= 0.163 < 0.20$.
나머지 세 쌍 모두 유의미한 단조 추세를 보인다($p < 0.01$).
이는 $g_{\min}$이 진폭 반상관의 \emph{정준 간격 메트릭}임을 확인한다.
\end{observation}

메커니즘은 투명하다: 하한
$A \ge 2/\Delta_{\min}^2$ (\cref{thm:lower_bound})은 순전히 대수적이며
높이-독립적이다. 하한이 예리하므로
($2H_1^{\mathrm{NN}}/A \approx 87\%$), 순위 상관은
기계적으로 하한에 의해 구동되어 높이-독립성을 물려받는다.
반면 $g_{\mathrm{right}}$ 메트릭은
감마 인자 오염($\sim \log(t)$ 가산적 이동)을 받아
높이에 따라 단조롭게 저하된다.

\begin{remark}[Cauchy 방법 검증]
$T \le 300$에서 $A_\La$의 독립적 Cauchy 윤곽 계산
($n = 136$, 반지름 $= 0.01$, 128점 구적,
\texttt{dps}$= 60$)은
$\rho(A_\La, g_{\mathrm{right}}) = -0.580$($p = 1.4 \times 10^{-13}$)을 주어,
원래 측정값 $-0.590$을 $|\Delta\rho| = 0.010$ 이내로 재현한다.
이 낮은-$t$ 값($-0.58$)과 전체 범위 값($-0.46$)의 차이는
위에서 기록된 높이 감쇠로 완전히 설명된다.
\end{remark}


% ============================================================================
\section{GUE 예측}
\label{sec:gue}
% ============================================================================

진폭--간격 상관의 랜덤 행렬 내용을 정량화하기 위해,
$300 \times 300$ GUE 행렬의 $1000$회 시뮬레이션을 수행하고,
$88{,}000$개의 bulk 고유값(중앙 $30\%$)을 추출하여
Hadamard 진폭
$A = \sum_{k \ne n} 1/(E_n - E_k)^2$을 계산한다
(최근접 이웃 $K \ge 5$개에서 수렴).

\begin{center}
\begin{tabular}{lccc}
\toprule
\textbf{측도} & \textbf{GUE 예측} & \textbf{관측} & \textbf{일치도} \\
\midrule
$\rho(A, g_{\mathrm{right}})$ & $-0.499 \pm 0.001$ & $-0.578$ & $86\%$ \\
$\rho(A_n, A_{n+1})$ & $+0.363$ & $+0.28$--$+0.42$ & 정확 \\
NN 기여 $2H_1^{\mathrm{NN}}/A$ & $68\%$ & $78\%$ & 방향적 \\
\bottomrule
\end{tabular}
\end{center}

GUE는 관측된 $g_{\mathrm{right}}$ 상관의 $86\%$를 재현한다.
잔차 $\Delta\rho \approx 0.08$의 두 후보 원인:

\begin{enumerate}
  \item \textbf{감마 인자 교차항.}\;
        감마 기여
        $\Delta A_\Gamma = -2 S_1^L \cdot \operatorname{Im}(\psi(\rho/2)/2)$는
        낮은 $t$에서 상관을 강화하는 높이-의존적 이동을 도입한다.
        직접 측정:
        $\rho(\Delta A_\Gamma, g_{\mathrm{right}}) = -0.70$
        ($p \approx 0$),
        교차항이 $14\%$ 잔차의 주요 원천임을 확인한다.
  \item \textbf{산술적 내용.}\;
        $S_1^2/A$ 비율이 GUE bulk($11\%$)와 $L$-함수(${\sim}23\%$)에서 다르며,
        이는 교대 합 $B$에서 산술 계수의 비소거를 반영한다.
\end{enumerate}


% ============================================================================
\section{편상관}
\label{sec:partial}
% ============================================================================

``공유 간격'' 효과를 내재적 진폭--진폭 결합에서 분리하기 위해,
$n$번째와 $(n\!+\!1)$번째 영점 사이의 간격을 통제한
편상관 $\rho(A_n, A_{n+1} \mid g_{n,n+1})$을 계산한다.

$\zeta(s)$($n = 910$)에 대해:
$\rho_{\mathrm{편상관}}(A_n, A_{n+1} \mid g_{n,n+1}) = +0.391$
($p = 1.3 \times 10^{-18}$).
공유 간격 효과를 제거한 후에도 인접 진폭들은 양의 상관을 유지한다.
메커니즘은 공유 최근접 이웃 항이다: $\gamma_n$과 $\gamma_{n+1}$ 모두
$\gamma_{n+2}$에 가깝다면, 두 진폭 모두 $\gamma_{n+2}$로부터 큰
$H_1$ 기여를 받는다.


% ============================================================================
\section{$\sigma$--$A$ 교차 검증}
\label{sec:crosscheck}
% ============================================================================

$\sigma$-프로파일 전개(\cite{paper2}의 따름정리 4.3)는
$A(\gamma)$에 대한 독립적 경로를 제공한다:
곡률 프로파일 $\kappa(\sigma, \gamma) = 1/(\sigma - \tfrac{1}{2})^2 + A + O((\sigma-\tfrac{1}{2})^2)$를
적합하면 Laurent 전개의 $A(\gamma)$와 같아야 하는 계수 $c(\gamma)$를 얻는다.

20개의 $\zeta(s)$ 영점에 대해:
$\rho(\text{$\sigma$-스캔 } c(\gamma),\; \text{Laurent } A(\gamma))
= 0.99999744$,
최대 상대 오차 $0.19\%$.
두 완전히 독립적인 계산 방법---하나는 $\sigma$-방향 적합,
다른 하나는 Cauchy 윤곽 적분---의 이 6자리 일치는
진폭 정의의 유효성을 검증하고 체계적 수치 인공물을 배제한다.


% ============================================================================
\section{논의}
\label{sec:discussion}
% ============================================================================

\subsection{결과 요약}

Hadamard 진폭 $A(\gamma)$가 자기동형 $L$-함수의 개별 영점의
의미 있는 불변량임을 확립했다:
\begin{enumerate}[label=(\roman*)]
  \item 무조건적 하한 $A \ge 2/\Delta_{\min}^2$
        (\cref{thm:lower_bound})은 Hadamard 곱 공식의 직접적 귀결이며
        미증명 가설을 필요로 하지 않는다.
  \item 반상관은 차수 1--4의 열 개의 $L$-함수에서 검증되어
        차수-독립 보편성을 확립한다.
  \item GUE 랜덤 행렬 앙상블이 상관의 $86\%$를 설명하며,
        잔차는 감마 인자 교차항에 기인함을 규명한다.
  \item 쌍 $(A_\La, g_{\min})$이 높이-안정한 반면
        $(A_\La, g_{\mathrm{right}})$은 감쇠하여,
        $g_{\min}$이 정준 메트릭임을 확인한다.
\end{enumerate}

\subsection{물리적 해석}

진폭 $A(\gamma)$는 영점의
\emph{스펙트럼 영향 반경}으로 해석될 수 있다:
$A$가 큰 영점은 국소 해석적 구조에 더 강한 영향을 미치며,
\cref{thm:lower_bound}은 이 영향이 최근접 이웃 거리의 제곱에
반비례함을 보여준다.
양자화 Hamiltonian의 Berry--Keating 스펙트럼 해석
$A \sim$ 국소 상태 밀도와의 유추는 시사적이나
아직 증명되지 않았다.

\subsection{열린 문제들}

\begin{enumerate}
  \item \textbf{$\rho \approx -0.57$의 해석적 도출.}\;
        Spearman 상관의 특정 값을 쌍 상관 함수로부터 도출할 수 있는가?
  \item \textbf{$2H_1/A$의 차수 스케일링.}\;
        비율이 $0.87$(GL(1))에서 ${\sim}0.67$(GL(4))로 감소하는데,
        이것이 $L$-함수의 차수로부터 예측될 수 있는가?
  \item \textbf{$\sigma$-국소화와 $A$--gap 이중성.}\;
        $\sigma$-프로파일 $\kappa(\sigma) = 1/\varepsilon^2 + A + \cdots$와
        간격 하한 $A \ge 2/\Delta_{\min}^2$은
        영점의 횡방향($\sigma$-방향)과 종방향($t$-방향) 구조 사이의
        이중성을 시사한다.
  \item \textbf{GL(5) 이상.}\;
        $\mathrm{sym}^4$(차수 5) 이상으로의 확장은
        Langlands 계층에서 보편성이 지속되는지 검증할 것이다.
\end{enumerate}


% ============================================================================
% 참고문헌
% ============================================================================
\begin{thebibliography}{99}

\bibitem{montgomery1973}
H.~L.~Montgomery,
``The pair correlation of zeros of the zeta function,''
\textit{Proc.\ Symp.\ Pure Math.}, vol.~24, pp.~181--193, 1973.

\bibitem{odlyzko1987}
A.~M.~Odlyzko,
``On the distribution of spacings between zeros of the zeta function,''
\textit{Math.\ Comp.}, vol.~48, pp.~273--308, 1987.

\bibitem{paper1}
Y.-H.~Kim,
``위상 양자화, 게이지 ODE와 $L$-함수 영점을 위한 $\xi$-다발 프레임워크 I,''
2026 (프리프린트).

\bibitem{paper2}
Y.-H.~Kim,
``위상 양자화, 게이지 ODE와 $L$-함수 영점을 위한 $\xi$-다발 프레임워크 II:
이론적 기초,''
2026 (프리프린트).

\bibitem{paper3}
Y.-H.~Kim,
``위상 양자화, 게이지 ODE와 $L$-함수 영점을 위한 $\xi$-다발 프레임워크 III:
확장과 보편성,''
2026 (프리프린트).

\bibitem{titchmarsh1986}
E.~C.~Titchmarsh,
\textit{The Theory of the Riemann Zeta-Function}, 2판
(D.~R.~Heath-Brown 개정). Oxford University Press, 1986.

\bibitem{iwaniec2004}
H.~Iwaniec, E.~Kowalski,
\textit{Analytic Number Theory}.
AMS Colloquium Publications, vol.~53, 2004.

\bibitem{mehta2004}
M.~L.~Mehta,
\textit{Random Matrices}, 3판. Academic Press, 2004.

\end{thebibliography}

\end{document}
